\section{Exercício prático: aplicação com persistência de dados}

Neste exercício, o objetivo é criar um Dockerfile para uma aplicação Node.js simples que recebe uploads de arquivos, gerar a imagem e garantir que os arquivos enviados pelo usuário não sumam ao reiniciar o container, utilizando volumes.

\subsection{Estrutura da aplicação}

Considere uma aplicação Node.js com Express que salva uploads em um diretório \texttt{/app/uploads}:

\begin{verbatim}
minha-app/
  Dockerfile
  package.json
  server.js
  uploads/       # diretório de uploads (persistido via volume)
\end{verbatim}

\subsection{Dockerfile}

\begin{verbatim}
FROM node:22-alpine
WORKDIR /app
COPY package.json ./
RUN npm install --production
COPY server.js ./
RUN mkdir -p /app/uploads
EXPOSE 3000
CMD ["node", "server.js"]
\end{verbatim}

\subsection{Build da imagem}

\begin{verbatim}
docker build -t minha-app .
\end{verbatim}

\subsection{Executar SEM volume (dados são perdidos)}

\begin{verbatim}
docker run -d -p 3000:3000 --name app-sem-volume minha-app
# ... enviar arquivos para a aplicação ...
docker rm -f app-sem-volume
# Os arquivos enviados são perdidos!
\end{verbatim}

\subsection{Executar COM volume (dados persistem)}

\begin{verbatim}
# Criar o volume
docker volume create uploads-data

# Iniciar o container com o volume montado em /app/uploads
docker run -d -p 3000:3000 --name app-com-volume \
    -v uploads-data:/app/uploads minha-app

# ... enviar arquivos para a aplicação ...

# Parar e remover o container
docker rm -f app-com-volume

# Iniciar novo container com o MESMO volume
docker run -d -p 3000:3000 --name app-restaurado \
    -v uploads-data:/app/uploads minha-app

# Os arquivos enviados ainda estão lá!
\end{verbatim}

\subsection{Conclusão}

A chave é identificar quais diretórios dentro do container contêm dados que devem sobreviver ao ciclo de vida do container (como uploads, bancos de dados, logs) e montar volumes nesses diretórios. Assim, mesmo que o container seja parado, removido ou recriado, os dados permanecem intactos.
